\section{Introduction}\label{sec:introduction}

Let $G$ be a finite group and let $p$ be a prime. The poset $\Ap(G)$
consists of its nonidentity elementary abelian $p$-subgroups, ordered by
inclusion. Thus each vertex is a subgroup isomorphic to $(\mathbb Z/p\mathbb
Z)^r$ for some $r\geq1$, and a simplex of its order complex is a strictly
increasing chain of such subgroups. Write $\Op(G)$ for the largest normal
$p$-subgroup of $G$.

Quillen's conjecture asserts that contractibility of this order complex
forces $\Op(G)\ne1$. It asks whether the topology of the elementary
subgroups detects the existence of a normal $p$-subgroup. We prove a
stronger statement.

\begin{theorem}\label{thm:main}
For every finite group $G$ and every prime $p$, if $\Op(G)=1$, then
\[
\widetilde H_*(\Ap(G);\Q)\ne0.
\]
Here reduced homology is augmented, so that the empty poset has
$\widetilde H_{-1}(\varnothing;\Q)=\Q$.
\end{theorem}

A contractible nonempty complex has zero reduced homology, and the empty
complex is not contractible. Theorem~\ref{thm:main} therefore gives a
positive resolution of Quillen's conjecture, without a restriction on
the prime or the composition factors of $G$.

The rational conclusion also determines acyclicity with other
coefficients. Since the order complex is finite, a nonzero rational
homology group comes from an integral homology group of positive free
rank. The universal coefficient theorem therefore gives nonzero reduced
homology over every field in the same degree; for the empty poset this
follows directly from augmentation. Conversely, a nontrivial $p$-core
makes $\Ap(G)$ contractible \citep[Proposition~2.4]{Quillen1978}.
Thus acyclicity over any fixed field characterizes the existence of a
nontrivial normal $p$-subgroup.

\paragraph{Background.}
Brown's work on Euler characteristics introduced the poset
$\mathcal S_p(G)$ of all nonidentity $p$-subgroups and related its Euler
characteristic to the order of a Sylow subgroup \citep{Brown1975}.
Quillen proved that the inclusion $\Ap(G)\subseteq\mathcal S_p(G)$ is a
homotopy equivalence \citep[Proposition~2.1]{Quillen1978}, so the two
posets carry the same homological information. He formulated the
contractibility conjecture in 1978 and established it for solvable
groups, groups of elementary abelian $p$-rank at most two, and groups
of Lie type in defining characteristic
\citep[Conjecture~2.9, Proposition~2.10, Theorem~3.1, and
Corollary~12.2]{Quillen1978}. For solvable groups with trivial $p$-core,
his proof produces nonzero homology in the largest possible degree.
This stronger conclusion, subsequently called the \emph{Quillen
dimension property}, became a central input to the component reductions.

The almost-simple case was established by
\citet{AschbacherKleidman1990}. Passing from simple components to an
arbitrary finite group requires control of their extensions and of the
homology that survives in the ambient group.
\citet{AschbacherSmith1993} developed this approach for $p>5$, reducing
the remaining problem to the dimension property for specified
elementary extensions of unitary groups.
\citet[Theorems~1.1 and~1.4]{PitermanSmith2025} extended the reduction
to all odd primes. \citet[Theorems~B and~C]{DiazRamos2026} proved the
required unitary result and deduced the rational homological assertion
for every finite group at odd primes. We retain a direct frame proof
of the required unitary dimension property. Its local statement,
Proposition~\ref{prop:unitary-odd}, covers the indicated almost-simple
groups with elementary abelian quotient, without assuming a splitting.

At two, successive component eliminations narrowed the remaining
families. \citet[Corollary~1.3]{PitermanSmith2022} showed that every
component of a minimal counterexample must admit an outer automorphism
of order two induced by the ambient group.
\citet[Theorem~1.1 and Corollary~1.2]{Piterman2024} established the
dimension property for elementary $2$-extensions of exceptional
Lie-type groups in odd characteristic, apart from a finite list in
characteristic three. In a minimal counterexample,
\citet[Theorem~A]{Piterman2026} then eliminated the remaining Lie-type
components in characteristic two. Combining these results with the
preceding reductions gives the classical components in characteristic
at least five listed in \citet[Theorem~B]{Piterman2026}.
We use this reduction in the precise form recalled in
Section~\ref{sec:reductions} and prove the remaining cases by our cycle
and propagation constructions.

Several of the methods behind these reductions are also central to the
present proof. \citet{SegevWebb1994} developed exact-sequence methods
relating subgroup posets to a normal subgroup and its centralizers.
\citet[Sections~3--5]{Piterman2021} developed full-chain homology
propagation, maximal faithful elementary extensions, and deletion of
overgroups, building on the propagation arguments of
\citet{AschbacherSmith1993}. Our propagation proof spells out the
coefficient and maximality conditions used by the constructions below.
The constructible cycles of \citet{DiazRamos2018} and the admissible
collections of \citet{DiazRamosMazza2022} likewise organize subgroup
flags by separating nonzero simplex coefficients from vanishing
boundary coefficients. The frame and decoration constructions here
use this chain-level viewpoint, with quantitative estimates that are
uniform as the dimension grows over a fixed field.

\paragraph{The cycle construction.}
A frame is a decomposition of a finite-dimensional vector space into
lines; when a form is present, the lines are nondegenerate and mutually
orthogonal. Independent scalar signs on the lines give an elementary
abelian projective subgroup. Its full subgroup flags support familiar
apartment chains. We impose relations between frames that differ only
inside a plane, so that the boundaries of their apartment chains cancel.

The main quantitative ingredient is a lower bound on the dimension of
the resulting coefficient space, uniform as the dimension of the
ambient space increases while the field stays fixed. This bound gives
room to impose further equations: adjoining commuting elementary
automorphisms creates new boundary faces, and their cancellation is a
linear condition on the frame coefficients. Counting the possible
completions of each partial choice bounds the number of such conditions.
The comparison produces a nonzero cycle on full subgroup flags, rather
than merely a nonzero collection of parameters describing apartments.

To obtain the dimension bound, we encode merger trees of frame lines
by multilinear Lie superbrackets. A Poincar\'e--Birkhoff--Witt argument
and a bound for the actual kernel of the generator--relation complex
turn elementary plane counts into the estimate of
Theorem~\ref{thm:frame-bound}. The supporting subspace and its isometry
type are retained throughout, so the same estimate can be used after
the descent imposed by specified commuting automorphisms.
In the symmetric case with field parameter at least five, the retained
fraction is at least $0.68^h$ in dimension $h$.

The geometric and algebraic ingredients have established precedents.
The unitary decomposition poset is studied by
\citet{PitermanWelker2024}; the relation between merger trees and
multilinear Lie representations appears in
\citet{Stanley1982,Robinson2004}. We identify the precise coefficient
space needed here and retain the supporting-subspace grading in the
PBW calculation. The generator--relation series method is related to
that of \citet{GolodShafarevich1964}; the additional kernel inequality
is proved in Section~\ref{sec:frames} before the logarithmic comparison
is used.

\paragraph{From a coefficient to homology.}
The constructions may first give a cycle in an automorphism group
larger than the chosen extension $H$ of a simple component $L$. Intersecting
arbitrary subgroup flags with $H$ could collapse them or cancel their
coefficients. Lemma~\ref{lem:restriction} avoids this by selecting a
supported elementary group with largest intersection index, taking a
residue at a fixed complement, and then intersecting the remaining
flag with $H$. On that residue the operation is injective, so one
nonzero full-flag coefficient survives.

At the prime two we work in a minimal counterexample and choose the
last subgroup $A$ of maximal rank in a family of elementary subgroups
generating faithful component extensions. Proper-subgroup induction supplies a
nonbounding cycle in the centralizer of $H=LA$. Its product with the
constructed cycle has a residue at the chosen flag that detects the
centralizer class. The deletion argument in
Lemma~\ref{lem:propagation} ensures that this residue still detects
nonbounding in the ambient group. This proves nonvanishing without
requiring $A$ to have the largest possible elementary rank in $G$.

For odd primes the cited reduction requires a different conclusion:
nonzero homology in the largest possible degree for certain unitary
extensions. We establish the needed rank formula before constructing
their cycles. This distinguishes the top-degree argument from the
propagation argument used at two.

Two parts of the proof take shorter or different routes through these
steps. The symplectic cycles are constructed directly in $H$ and need
no restriction from a larger group. In the last branch for
$\PSL_4(5)$, an equivalent-poset construction instead retains the join
of the component and centralizer homology. Section~\ref{sec:exceptional}
states the precise conditions that lead to this branch.

\paragraph{Organization.}
Section~\ref{sec:reductions} fixes conventions and states the reduction
theorems. Section~\ref{sec:propagation} proves the propagation lemma.
Sections~\ref{sec:frames} and~\ref{sec:chains} develop frame cycles,
their dimension bounds, and the decoration and restriction procedures.
Sections~\ref{sec:linear}--\ref{sec:orthogonal} establish the required
linear, unitary, symplectic, and orthogonal cases.
Section~\ref{sec:exceptional} treats the remaining small orthogonal
case, realized as $\PSL_4(5)$, and Section~\ref{sec:conclusion} assembles
the proof of Theorem~\ref{thm:main}.
